\documentclass[12pt]{article}
\usepackage[T1]{fontenc}
\usepackage[utf8]{inputenc}
\usepackage[brazil]{babel}
\usepackage[margin=1in]{geometry}
\usepackage{ragged2e}
\usepackage[normalem]{ulem}
\setlength{\parindent}{0pt}
\setlength{\parskip}{0.8\baselineskip}
\begin{document}
\justifying
\noindent Verifica-se que a questão objeto da controvérsia jurídica suscitada no recurso manejado pela parte recorrente esteve afetada no representativo da controvérsia - \textbf{Tema n. }\textbf{244} da TNU - (PU no processo n. PEDILEF 5002880-91.2016.4.04.7105/RS),  afetado em 06/11/2019 foi julgado (trânsito em julgado em RE 1413882/RS), por meio do qual foi elucidada a seguinte questão:

\noindent \textit{"Saber se o auxílio-alimentação, pago em espécie e com habitualidade, por meio de vale-alimentação ou na forma de tickets, tem natureza salarial e integra o salário de contribuição para fins de cálculo da renda mensal inicial (RMI)."}

\noindent Restou firmada a seguinte tese:

\noindent \textit{"I) Anteriormente à vigência da Lei n. 13.416/2017, o auxílio-alimentação, pago em espécie e com habitualidade ou por meio de vale-alimentação/cartão ou tíquete-refeição/alimentação ou equivalente, integra a remuneração, constitui base de incidência da contribuição previdenciária patronal e do segurado, refletindo no cálculo da renda mensal inicial do benefício, esteja a empresa inscrita ou não no Programa de Alimentação do Trabalhador - PAT; II) A partir de 11/11/2017, com a vigência da Lei n. 13.416/2017, que conferiu nova redação ao § 2º do art. 457 da CLT, somente o pagamento do auxílio-alimentação em dinheiro integra a remuneração, constitui base de incidência da contribuição previdenciária patronal e do segurado, refletindo no cálculo da renda mensal inicial do benefício, esteja a empresa inscrita ou não no Programa de Alimentação do Trabalhador - PAT." (Tema 244 TNU).}

\noindent Com efeito, depreende-se da leitura do voto condutor do acórdão que o mesmo se encontra em conformidade com o decidido no  tema pela TNU.

\noindent A proposito, consta do acordao hostilizado: (...) RMI - Renda Mensal Inicial, RMI - Renda Mensal Inicial, Reajustes e Revisões Específicas, DIREITO PREVIDENCIÁRIO, Competência, Jurisdição e Competência, DIREITO PROCESSUAL CIVIL E DO TRABALHO

\noindent Nessas condições, o conteúdo do Acórdão recorrido é consentâneo com o entendimento da TNU, firmado em julgamento de recurso representativo de controvérsia, o que atrai a incidência da regra prevista pelo art. 14, III, \textit{b,} da Resolução CJF n. 586/2019:

\noindent \textit{\textbf{Art. 14.}}\textit{ Decorrido o prazo para contrarrazões, os autos serão conclusos ao magistrado responsável pelo exame preliminar de admissibilidade, que deverá, de forma sucessiva: (...) }\textit{\textbf{III }}\textit{- negar seguimento a pedido de uniformização de interpretação de lei federal interposto contra acórdão que esteja em conformidade com entendimento consolidado: }\textit{\textbf{b)}}\textit{ em }\uline{\textit{\textbf{recurso representativo de controvérsia pela Turma Nacional de Uniformização}}}\textit{ ou em pedido de uniformização de interpretação de lei dirigido ao Superior Tribunal de Justiça;(...).}

\noindent Posto isso, com arrimo no \uline{\textbf{art. 14, III, }}\uline{\textit{\textbf{b}}}\textit{,} da Resolução CJF n. 586/2019 c/c art. 1.030, I, do CPC, \textbf{NEGO SEGUIMENTO }\textbf{a}o\textbf{ PU}\textbf{.}

\noindent Intimem-se. Aguarde-se o decurso do prazo recursal. Após, certifique-se o trânsito em julgado e devolva-se ao juízo de origem.
\end{document}
